\section*{Konklusjon}
\addcontentsline{toc}{section}{Konklusjon}

Kvile AS er en serie-/batchprodusent med høy variantbredde og tydelig sesongvariasjon. Dagens MTO-styring for Vegvesenet er hensiktsmessig og bør videreføres. Hovedanbefalingen er å flytte butikkmarkedets KODP fra ferdigvarelager til komponentlager gjennom ATO. Dette reduserer antall lagervarer med ca.\ 76\,\% og angriper hovedutfordringen: mange produktvarianter, sesongvariasjon og kapitalbinding i ferdigvarer.

Vi vurderer fabrikklayouten som primært funksjonell, med linjeorganiserte lakkanlegg. Ved overgang til ATO bør komponentlager plasseres nær endemontering og pakking. Bærekraften kan styrkes gjennom returemballasje, resirkulert aluminium, bedre ruteplanlegging og modulært design for reparasjon. Flere av tiltakene støtter også Lean-prinsipper ved å redusere lager, transport og materialavfall.
